\section{Conclusion}
% ══════════════════════════════════════════════════════════════

We studied LLMs as implicit factor-like priors for cross-sectional stock ranking, bridging LLM capability evaluation and factor diagnostics.
The central lesson is that model choice is an asset-pricing choice: switching the LLM changes not only ranking accuracy, but also the growth, momentum, value, profitability, and investment exposures being tested.
Across 11 included models, expected-return prompting is the primary lens: o3 reaches a 466.0\% net long--short CAGR, Sharpe 5.54, and IC 0.594 after charging both long- and short-leg turnover, while GPT-4o and GPT-4o mini produce weak or negative signals.
The input ablation shows that market summaries carry most of the pricing signal and factor-profile coherence, whereas fundamentals-only prompts weaken both alpha and style structure.
For finance benchmarks, the implication is to report ordering quality, investable spread performance, and factor diagnostics together rather than treating an LLM ranking as a black-box forecast.
Future work should extend to broader rolling universes, live model refreshes, and disclosure-text information axes while preserving the same audit trail from prompt to ranking, portfolio, and factor exposure.

% ══════════════════════════════════════════════════════════════
